\documentclass[11pt,a4paper]{article}
\usepackage[margin=23mm]{geometry}
\usepackage[T1]{fontenc}
\usepackage[utf8]{inputenc}
\usepackage{lmodern,microtype,amsmath,amssymb,booktabs,tabularx,array}
\usepackage{graphicx,xcolor,listings,enumitem}
\usepackage[section]{placeins}
\usepackage[colorlinks=true,linkcolor=blue!50!black,urlcolor=blue!50!black]{hyperref}
\usepackage{fancyhdr}
\pagestyle{fancy}
\fancyhf{}
\lhead{Statistical learning: worked notebook explanations}
\rfoot{\thepage}
\setlength{\headheight}{14pt}
\setlength{\parindent}{0pt}
\setlength{\parskip}{6pt}
\setlength{\emergencystretch}{2em}
\setlist{itemsep=3pt,topsep=4pt}
\lstset{language=Python,basicstyle=\small\ttfamily,breaklines=true,
 columns=fullflexible,keepspaces=true,showstringspaces=false,
 frame=single,rulecolor=\color{gray!35},backgroundcolor=\color{gray!4},
 commentstyle=\color{green!35!black},keywordstyle=\color{blue!65!black},
 aboveskip=8pt,belowskip=8pt}
\newcommand{\code}[1]{\texttt{\detokenize{#1}}}
\newcommand{\RSS}{\operatorname{RSS}}
\newcommand{\RMSE}{\operatorname{RMSE}}
\newcommand{\MAE}{\operatorname{MAE}}
\newcommand{\E}{\mathbb{E}}
\newcommand{\ind}[1]{\mathbf{1}\{#1\}}
\newcommand{\fig}[2]{\begin{figure}[htbp]\centering
\includegraphics[width=\linewidth,height=0.39\textheight,keepaspectratio]{figures/#1}
\caption{#2}\end{figure}}
\newcommand{\notebooktitle}[3]{\begin{center}{\LARGE\bfseries #1\par}\vspace{5pt}
{\large Problem statement and worked resolution\par}\vspace{5pt}
Companion to \code{#2}\par\end{center}
\textbf{How to use this explanation.} #3 Section numbers follow the notebook's numbered sections where present. Numerical values are rounded and refer to the recorded run, not universal performance claims.}
\newcommand{\methods}[1]{\subsection*{Main Python calls}#1}

\begin{document}
\notebooktitle{4B: PCA compression and denoising}{04_B_pca_compression.ipynb}{This explanation connects the numerical identities to projection geometry and derives why reconstruction of observations and recovery of signal can favor different dimensions.}
\section*{Problem statement}
We first reproduce an exact two-feature PCA calculation. We then generate a three-dimensional signal embedded in 20 features,
add isotropic noise, and compare several fixed rules for how many components to keep.
The clean signal is known only because the data are simulated. A rule consulting clean validation rows is called an oracle diagnostic;
it is not presented as a procedure available for ordinary noisy data.

Let X have n rows and p columns. Centering gives $Z=X-\mathbf1\bar x^T$, with fitting mean $\bar x$.
Principal component analysis (PCA) finds unit directions of large sample variance. The direction matrix $V_q$ has shape $(p,q)$,
the score matrix $T_q=ZV_q$ has shape $(n,q)$, and reconstruction is
\[\widehat X_q=\mathbf1\bar x^T+T_qV_q^T.\]
The squared Frobenius norm sums entrywise squared errors. Mean squared error (MSE) in this notebook divides that total by np.

\section{The exact calculation and why it verifies more than signs}
The four rows are $(3,1),(1,3),(-1,-3),(-3,-1)$, whose mean is zero.
Their sample covariance is
\[S=\frac1{3}\begin{pmatrix}20&12\\12&20\end{pmatrix}.\]
Multiplying by $(1,1)^T$ gives $(32/3)(1,1)^T$; multiplying by $(1,-1)^T$ gives $(8/3)(1,-1)^T$.
After normalization these are orthonormal directions. The first proportion of variance explained (PVE) is
$32/(32+8)=0.8$, and its scores are $2\sqrt2,2\sqrt2,-2\sqrt2,-2\sqrt2$ under the displayed sign.
Projecting back gives $(2,2),(2,2),(-2,-2),(-2,-2)$.

For each row, the residual has squared length 2. Thus total error is 8, error per row is 2,
and error per entry is 1. These are three different quantities obtained from the same residual matrix.
The singular value decomposition (SVD) has singular values $\sqrt{32}$ and $\sqrt8$.
Its returned third array in NumPy is $V^T$, so slicing its first row selects a direction, not a column of scores.

A direction and its negative define the same line. A strong numerical check compares
$V_qV_q^T$, the projection matrix, rather than signs. This also handles rotations among components with repeated eigenvalues.
The notebook verifies the projector, reconstruction, eigenvalues, PVE, and error identity. A matching scatterplot alone would not detect every transpose or scaling mistake.

\methods{\code{PCA.fit} learns \code{mean_} and direction rows in \code{components_}.
\code{transform} subtracts that learned mean and returns scores; \code{inverse_transform} maps scores back.
\code{svd_solver="full"} makes the dense decomposition choice explicit, and whitening remains off.}

\section{What signal and noise mean in the simulation}
A thin QR decomposition of a seeded $20\times3$ Gaussian matrix creates B with $B^TB=I_3$.
Draw latent coordinates with variances 9, 4, and 1. The clean signal row is $s=Ba$ and the observed row is
$x=s+\epsilon$, where $\epsilon$ has covariance $0.25I_{20}$ and is independent of a.
Therefore population covariance is
\[\operatorname{Cov}(x)=B\operatorname{diag}(9,4,1)B^T+0.25I_{20}.\]
Its eigenvalues are $9.25,4.25,1.25$, followed by seventeen values of 0.25.
Total variance is $14+20(0.25)=19$. Even a perfectly known signal subspace retains only
$14.75/19\approx77.63\%$ of observed variance. The remaining variance is noise, not necessarily information worth preserving.
This population calculation anticipates why a 90\% compression cutoff keeps more than three directions.

The split contains 360 fitting rows, 120 validation rows, and 120 test rows. The PCA mean and full 20-direction basis
are learned only on fitting observations. Clean rows do not enter the PCA fit. The fitting mean norm is about 0.1091;
it need not be exactly the population mean zero in a finite sample.

No standardization is used for the primary fit because the simulation specifies equal noise variance in the original coordinate system.
Standardization would weight squared errors differently. ``Always scale before PCA'' is therefore too broad a rule:
first identify the geometry and noise model relevant to the task.

\section{Why the two error curves behave differently}
For each $q=0,\ldots,20$, the notebook uses the same learned mean and nested basis.
At q=0, multiplication through an empty direction matrix yields the fitting mean for every row.
At q=20, the directions span all 20 feature coordinates, so every input row is reconstructed exactly, apart from rounding.
This remains true for new rows; no test fitting is required.

For fitting rows, orthogonality gives
\[
\|Z-ZV_qV_q^T\|_F^2=(n_{\rm train}-1)\sum_{j>q}\lambda_j.
\]
The factor is 359 here. The fitting eigenvalues do not determine the exact residual sum on a different validation sample.
However, for any fixed row and nested orthogonal projections, adding a direction cannot increase error to that same observed row.
The validation noisy-error curve is therefore also nonincreasing, without invoking any learning-performance guarantee.

Why can signal error increase? If P is a fixed rank-q projector, independent of new isotropic noise, and the true mean is known,
\[
Px-s=-(I-P)s+P\epsilon.
\]
The two parts are orthogonal, and $\mathbb E\|P\epsilon\|^2=\sigma^2\operatorname{tr}(P)=q\sigma^2$.
Thus
\[\mathbb E\|Px-s\|^2=\|(I-P)s\|^2+q\sigma^2.\]
Dropping a signal direction creates bias; retaining an extra noise direction adds variance.
With the exact rank-three projector, ideal per-entry signal error is $3(0.25)/20=0.0375$.
Estimated directions and means add further finite-sample effects, so this is a reference calculation rather than a prediction of every run's precise MSE.

\fig{04_B_pca_compression_1.png}{Cumulative explained variance and reconstruction curves. More components improve fidelity to observed noise; the clean-signal curve has a different objective.}
The recorded validation values are:
\begin{center}\begin{tabular}{rrr}\toprule
q&Noisy-observation MSE&Clean-signal MSE\\\midrule
0&1.030729&0.776397\\3&0.224575&0.036037\\
11&0.128518&0.129609\\20&approximately 0&0.257001\\\bottomrule\end{tabular}\end{center}
The fitting CPVE rule selects q=11. The clean-validation oracle selects q=3, agreeing with the known simulation dimension.
The same agreement need not hold under a different realization or noise structure.

\section{Read the final test as an assessment of frozen policies}
\begin{center}\begin{tabular}{lrrr}\toprule
Policy&q&Noisy MSE&Signal MSE\\\midrule
Mean only&0&0.998555&0.760701\\
Known signal dimension&3&0.224424&0.038227\\
90\% fitting variance&11&0.118193&0.143124\\
Clean-validation oracle&3&0.224424&0.038227\\
Full reconstruction&20&approximately 0&0.259112\\\bottomrule\end{tabular}\end{center}
The 90\% policy has lower error to observations than q=3 but considerably worse signal recovery.
Full reconstruction merely reproduces the observed test noise; its signal MSE equals the measured noise MSE, about 0.2591.
This is close to the generating variance 0.25 but need not equal it exactly.
The two q=3 rows are identical because they select the same frozen representation, not because one is an independently fitted replicate.

The Frobenius distance between the learned rank-three projector and $BB^T$ is about 0.15197.
Both projectors have rank 3; the comparison describes orientation error without arbitrary basis signs.
For rank-q projectors P and Q,
$\|P-Q\|_F^2=2q-2\operatorname{tr}(PQ)$; smaller distance means greater overlap of their subspaces.
This diagnostic is available because B is known in the simulation, not because PCA ordinarily reveals a true hidden basis.

\fig{04_B_pca_compression_2.png}{One held-out row. Connected feature indices aid comparison; they do not turn the feature vector into a time series.}
For 600 stored rows, q=3 uses $600(3)+20(3)+20=1880$ numerical entries versus $600(20)=12000$.
This is about 15.7\% of the original entry count, including direction and mean overhead.
The clean basis B is not needed to decode a reconstruction; the learned basis and mean are.
Compressed file size also depends on numerical precision and format overhead, which this count does not measure.

\section{Changing units changes an unscaled optimization}
Multiplying feature 1 by 100 changes its squared-error weight by a factor of 10,000.
The first raw PVE rises from about 0.481773 to 0.998630, showing how a unit convention can dominate the leading direction.
For a positive multiplier a, the fitted mean and standard deviation become $a\bar x$ and $as$,
so $(ax-a\bar x)/(as)=(x-\bar x)/s$. Standardized inputs and their PCA projectors consequently agree after a corresponding refit.
The standardized first PVE remains about 0.370997 in both unit systems.
This invariance does not demonstrate that standardized PCA is the best denoiser for the original isotropic-noise model.

\section*{Documentation and practical advice}
The \href{https://scikit-learn.org/1.8/modules/generated/sklearn.decomposition.PCA.html}{PCA reference}
distinguishes singular values, sample-variance eigenvalues, variance ratios, and component rows.
The \href{https://scikit-learn.org/1.8/modules/generated/sklearn.preprocessing.StandardScaler.html}{StandardScaler reference}
uses a divisor-n variance convention; it is separate from PCA's divisor-$n-1$ covariance convention.
For complete nonconstant columns this changes all standardized coordinates by a common factor compared with sample scaling,
so PVE and direction orientation agree, while absolute score variances do not.

Preserve the mean and basis together; do not refit them on a prediction batch.
Keep whitening off when checking the simple score formula $ZV_q$.
A very high explained variance ratio is not a guarantee of downstream prediction quality.
Finally, noisy validation reconstruction alone cannot select a denoising dimension here: its optimum is the full reconstruction.
A usable denoising selection method needs additional assumptions or information, which the notebook deliberately does not pretend to supply.
\end{document}
